% Transcription — fonds Alexandre Grothendieck, shelfmark 125, batch 1 (pages 1-8).
% Established from the facsimile archives/batches/125/batch-01.pdf (a mirror of
% https://grothendieck.umontpellier.fr/125.pdf), per the skill
% transcribe-grothendieck. The folder is eight pages and this is its only batch.
%
% The cover (page 1) reads "Dualité locale et dualité projective", with one word
% blotted out between "dualité" and "projective". The inventory records the
% folder as "Dualité projective" alone, so the cover carries a title one half
% longer than the catalogue's, and the folder's own words are what open this
% file. The cover gets no \page{}: the gap in the numbering is the record.
%
% Six leaves follow, and they are two documents interleaved. Pages 2, 4, 5, 7
% and 8 are his hand. Pages 3 and 6 are leaves of a typescript, headed
% "- III-18 -" and "- III-5 -", annotated in ink and pencil and marked with a
% yellow highlighter. They are NOT unrelated paper: III-5 proves
% H^p(U,F) = H^{p+1}((f),M) for U the complement of V(f), which is precisely
% the local-cohomology machinery the handwritten leaves are using. They are
% therefore transcribed, following folder 54, rather than skipped as the
% Bourbaki leaf of folder 44 was. Their typed prose is compressed to its
% numbered statements and displays, per the skill's rule on crowded pages;
% every mark made on them by hand is given in full.
%
% One thing the edition records and does not decide: the hand that annotates
% the two typescript leaves is round, careful and upright, and does not look
% like the fast angular hand of pages 2, 4, 5, 7 and 8. Whose it is, the folder
% does not say.
%
% Pages 7 and 8 are the piece most worth having: a programme in six numbered
% items setting out where formal schemes are needed — the two fundamental
% theorems of proper morphisms, the formal theory of Modules, the Lefschetz
% relation between a variety and its formal completion along a hyperplane
% section, the duality formalism, formal groups, and the Greenberg-Néron
% types. Item 2 is circled and item 3 arrowed, so the order on the page is not
% the order he settled on.
%
% Pass: Opus 5 (claude-opus-5), 2026-08-30 — first pass, unchecked against the pages by a human.
\documentclass[11pt,a4paper]{article}
% Le préambule commun à toutes les transcriptions du fonds Grothendieck.
%
% Il tient en une idée : **l'incertitude se compose, elle ne se résout pas.**
% Une transcription qui lisse un mot douteux a détruit la seule chose pour
% laquelle on la faisait — un lecteur doit pouvoir savoir, à chaque endroit,
% ce qui était sur le feuillet et ce qui a été ajouté. Les six macros qui
% suivent sont donc l'appareil critique, et non de la décoration.
%
% Les mêmes commandes sont reconnues par `scripts/render.mjs`, qui produit la
% vue de lecture du site. Une macro ajoutée ici doit l'être aussi là-bas,
% faute de quoi le rendu échouera bruyamment — ce qui est voulu : un rendu
% silencieusement faux est pire qu'un rendu absent.

\ProvidesPackage{grothendieck}[2026/08/08 Transcriptions du fonds Grothendieck]

\usepackage{amsmath,amssymb,amsthm}
\usepackage{mathrsfs}
\usepackage{stmaryrd}
% Les diagrammes commutatifs sont partout dans ce fonds : c'est la langue
% même de la géométrie algébrique de Grothendieck, et une transcription qui
% les rendrait en prose ne transcrirait rien.
\usepackage{tikz-cd}
% Français par défaut — c'est la langue des feuillets — mais l'anglais est
% chargé aussi : la traduction et le résumé sont des documents anglais, et la
% typographie française (espaces avant les deux-points) y serait une faute.
% Ces documents basculent par \selectlanguage{english} en tête de corps.
\usepackage[english,french]{babel}
\usepackage{fontspec}
\usepackage{geometry}
\geometry{a4paper,margin=25mm,marginparwidth=28mm}
\usepackage{marginnote}
\usepackage{xcolor}
% ulem, not soul. `soul`'s \st and \ul are built on a character-by-character
% scan that XeLaTeX's font machinery breaks: the document fails with an
% undefined control sequence rather than degrading. ulem does the same two jobs
% and is Unicode-engine safe. `normalem` stops it hijacking \emph, which the
% transcriptions use for Grothendieck's own emphasis.
\usepackage[normalem]{ulem}
\usepackage{hyperref}
\hypersetup{colorlinks=true,linkcolor=black,urlcolor=black,pdfborder={0 0 0}}

% --- Métadonnées du lot -----------------------------------------------------
% Reprises telles quelles par le script de rendu, qui les affiche en tête de
% la vue de lecture. Elles fixent ce que le fichier prétend couvrir : sans
% elles, une transcription flotte sans point d'attache dans un fonds de seize
% mille feuillets.

\newcommand{\folder}[1]{\gdef\grothFolder{#1}}
\newcommand{\batch}[1]{\gdef\grothBatch{#1}}
\newcommand{\leaves}[2]{\gdef\grothFirst{#1}\gdef\grothLast{#2}}
\newcommand{\foldertitle}[1]{\gdef\grothFoldertitle{#1}}
\gdef\grothFolder{?}\gdef\grothBatch{?}\gdef\grothFirst{?}\gdef\grothLast{?}\gdef\grothFoldertitle{}

% --- Le résumé --------------------------------------------------------------
%
% Il ouvre la lecture modernisée, sous le titre « Résumé », et il est composé
% en petit. Ce n'est pas un ornement : c'est le seul passage du document qui
% ne soit pas des mathématiques, et un lecteur doit pouvoir le reconnaître
% avant de l'avoir lu — puis le sauter s'il connaît déjà le sujet, ou s'y
% arrêter s'il ne le connaît pas. Le filet qui le suit marque la reprise du
% corps.

\newenvironment{resume}
  {\small\setlength{\parskip}{0.45em}}
  {\par\medskip\hrule\medskip}

% --- Mots-clés ---------------------------------------------------------------
%
% En anglais, en clôture du résumé de la lecture modernisée : le vocabulaire
% moderne sous lequel chercher ce que les feuillets construisent. Le manifeste
% du site (`npm run manifest`) les extrait pour étiqueter la cote entière —
% cette ligne est la source unique des tags, il n'y a pas de fichier à côté.
\newcommand{\keywords}[1]{\par\smallskip\noindent
  {\footnotesize\textsc{Keywords} --- \textit{#1}}\par}

% --- L'appareil critique ----------------------------------------------------

% Le numéro de feuillet, en marge. C'est l'ancre qui permet de régler un
% désaccord sur une formule en regardant *un* feuillet plutôt qu'une chemise
% de six cents. Le site s'en sert aussi pour tourner les pages du fac-similé
% quand on fait défiler la transcription.
\newcommand{\leaf}[1]{%
  \marginnote{\footnotesize\color{gray!70}\textsf{#1}}%
  \label{leaf:#1}\ignorespaces}

% Illisible : on ne devine pas. Un mot inventé qui se lit comme les autres est
% le pire résultat possible.
\newcommand{\ill}{\textcolor{red!60!black}{[\,\ldots\,]}}

% Lecture douteuse : le mot est donné, et signalé comme incertain.
\newcommand{\uncertain}[1]{\uline{#1}}

% Ajout de l'éditeur — une lettre manquante, un mot sous-entendu.
\newcommand{\add}[1]{\textcolor{blue!50!black}{[#1]}}

% Biffé par Grothendieck lui-même. Ce qu'il a rayé fait partie du document :
% une rature dit souvent où la pensée a changé de direction.
\newcommand{\struck}[1]{\sout{#1}}

% Note du transcripteur, sur l'état matériel du feuillet.
\newcommand{\note}[1]{\par\smallskip{\small\color{gray!60!black}\textit{[#1]}}\par\smallskip}

% Note marginale de Grothendieck, distincte du corps du texte.
\newcommand{\marginal}[1]{\par\smallskip\noindent\hspace{1em}%
  {\small\color{gray!60!black}#1}\par\smallskip}

% --- Titre ------------------------------------------------------------------

\AtBeginDocument{%
  \begin{center}
    {\large\bfseries Fonds Alexandre Grothendieck --- cote n\textsuperscript{o}~\grothFolder}\\[2pt]
    {\itshape\grothFoldertitle}\\[4pt]
    {\small feuillets \grothFirst--\grothLast\ (lot \grothBatch)}\\[2pt]
    {\footnotesize\color{gray!60!black}Transcription établie d'après le fac-similé numérisé
     par l'Université de Montpellier.\\
     Les lectures incertaines sont soulignées, les illisibles notées [\,\ldots\,],
     les ajouts entre crochets.}
  \end{center}
  \vspace{1em}}

\endinput


\folder{125}
\batch{1}
\pages{1}{8}
\dating{[avant 1970]}
\watermark{Édition de démonstration}
\foldertitle{Dualité projective : notes manuscrites (s.d.)}

\begin{document}

\section*{Dualité locale et dualité \struck{\ill{}} projective}

\note{titre de sa main, sur le premier feuillet, qui ne porte rien d'autre : le
feuillet est une chemise et ne reçoit donc pas de \texttt{page}. Un mot est
biffé d'un pâté d'encre entre \emph{dualité} et \emph{projective} et ne se
laisse pas lire. L'inventaire de Montpellier ne retient que
\guillemotleft{} Dualité projective \guillemotright{} : la chemise porte donc
un titre plus long que le catalogue}

\page{2}

\note{ce feuillet n'est pas une rédaction mais une carte : six énoncés encadrés,
reliés par des équivalences verticales, chacune annotée de ce qui la justifie.
On le lit de haut en bas dans la colonne de gauche, la colonne de droite
donnant à chaque étage la réduction au cas $F = \mathcal{O}(n)$, $n$ grand.
Les chapeaux désignent le complété formel le long de $Y$}

\[
  \mathbf{P}^{n} \supset X^{r} \supset Y^{m}, \qquad r > m > 0
\]

\[
  H^{p}_{Y}(X, F) = 0 \quad \text{pour } p > \ill{}
\]

$X$ de dimension $r$.

\[
  H^{p}(X^{r}, F) \longrightarrow H^{p}(X^{r} - Y, F) \longrightarrow
  H^{p+1}_{Y}(X^{r}, F)
\]

\note{la suite exacte de cohomologie locale, qu'il marque d'une astérisque : c'est
elle qui fait passer d'un étage à l'autre de tout ce qui suit}

Premier encadré :
\[
  H^{p}(X^{r} - Y, F) = 0 \quad \text{pour } p \geqslant m
\]
et, à droite, sa réduction :
\[
  H^{p}(X - Y, \mathcal{O}(n)) = 0 \quad \text{pour } n \text{ grand},\ p \geqslant m .
\]

Deuxième encadré :
\[
  H^{p}_{Y}(X^{r}, F) \longrightarrow H^{p}(X^{r}, F) \quad
  \begin{cases}
    \text{bijectif pour } p > m \\
    \text{surjectif pour } p = m
  \end{cases}
\]
et, à droite : \uncertain{réduit à} bijectif pour $F = \mathcal{O}(n)$,
$n$ grand.

\note{sous cet encadré de droite, accolé d'un trait :
\guillemotleft{} N.B. Le cas $p = m$ peut être omis, car
$H^{m}(X^{r}, \mathcal{O}(-n)) = 0$ pour $n$ grand \guillemotright}

Troisième encadré, relié au deuxième par une équivalence portant, entourée, la
mention \guillemotleft{} dualité \guillemotright{} :
\[
  \mathrm{Ext}^{q}_{\mathcal{O}_{\hat{X}^{r}}}
    \!\left(\hat{X}^{r}, \hat{F}, \hat{\Omega}^{r}_{X^{r}/k}\right)
  \;\xleftarrow{\ \sim\ }\;
  \mathrm{Ext}^{q}_{\mathcal{O}_{X^{r}}}
    \!\left(X^{r}, F, \Omega^{r}_{X^{r}/k}\right)
\]
\[
  \begin{cases}
    \text{bijectif pour } q < r - m \\
    \text{injectif pour } q = r - m
  \end{cases}
\]

\note{c'est le pivot du feuillet : l'annulation de la cohomologie de $X - Y$
en degrés $\geqslant m$ et la comparaison des Ext à valeurs dans
$\Omega^{r}$ le long du complété formel sont le même énoncé, lu des deux
côtés de la dualité}

Quatrième encadré :
\[
  H^{q}\!\left(X^{r}, \Omega(n)\right) \longrightarrow
  H^{q}\!\left(\hat{X}^{r}, \hat{\mathcal{O}}(n)\right) \quad
  \begin{cases}
    \text{bijectif pour } q < r - m \\
    \text{injectif pour } q = r - m
  \end{cases}
  \quad n \text{ grand}
\]

Cinquième encadré :
\[
  H^{q}(X^{r}, F) \longrightarrow H^{q}(\hat{X}^{r}, \hat{F}) \quad
  \begin{cases}
    \text{bij. pour } q < r - m \\
    \text{inj. pour } q = r - m
  \end{cases}
  \quad \text{pour } F \text{ loc. libre}
\]

\marginal{$H^{p}(X^{r}, \mathcal{O}(n)) \xrightarrow{\ \sim\ }
H^{p}(\hat{X}^{r}, \hat{\mathcal{O}}(n))$, et
$H^{i}(\hat{X}^{r}, \hat{\mathcal{O}}(n)) = 0$ pour $n$ grand, $0 < i < r-m$}

\note{note de sa main, entourée d'un trait, au bas du feuillet}

\marginal{N.B. \ill{} $H^{p}(X^{r} - Y, \Omega(n)) = 0$ si $0 < p < r-m$ \ill{}
$H^{p}(X^{r}, \Omega(n)) \longrightarrow H^{p}(X^{r} - Y, \Omega(n))$ \ill{}
$n < r-1$}

\note{cette note court en diagonale au bas du feuillet, sur cinq ou six lignes,
et n'est lisible que par fragments}

\ill{} $[\,m = \dim Y\,]$ \ill{}

\page{3}

\note{feuillet dactylographié, en tête \guillemotleft{} -- III-18 -- \guillemotright,
annoté à la main et souligné au feutre jaune. La prose tapée est ici resserrée
sur ses énoncés et ses formules ; les mentions manuscrites sont données en
entier. La main qui annote est ronde et posée, et ne ressemble pas à celle des
feuillets 2, 4, 5, 7 et 8 ; l'édition ne décide pas à qui elle est}

Le feuillet établit qu'une base $(d_{i} \otimes 1)_{1 \leqslant i \leqslant n}$
de $D \otimes_{C} C_{y}$ peut être prise formée de restrictions de sections de
$\mathcal{O}_{X}$ \emph{au-dessus de $X$}. Le préschéma $f^{-1}(y)$ a un espace
de base formé d'un nombre fini de points $x_{j}$
($1 \leqslant j \leqslant m$ ; II, 2, 7, prop.~17) formant un sous-espace
\emph{discret} de $f^{-1}(U)$ ; l'algèbre $D \otimes_{C} C_{y}$, qui est
l'anneau de ce schéma affine, est donc composée directe des sous-algèbres
$\mathcal{O}_{x_{j}}$ (I, 2, 7, prop.~19 et I, 4, 1, prop.~4).

\note{la prose intermédiaire, sur les $f_{j} \in \Gamma(X, \mathcal{O}_{X}) = B$
tels que $x_{j} \in X_{f_{j}} \subset f^{-1}(U)$, sur l'entier $q$ pour lequel
$d_{i} f_{j}^{q}$ se prolonge en une section $b_{i} \in B$ (I, 1, 4, th.~2) et
sur le fait que les $b_{i} \otimes 1$ forment encore une base, est ici
résumée ; une phrase du tapuscrit est biffée à la machine et une autre à la
main}

Cela étant, considérons l'homomorphisme
$\Gamma(Y, \mathcal{O}_{Y}^{n}) \to \Gamma(X, \mathcal{O}_{X}) = B$ faisant
correspondre $b_{i}$ au $i$-ème élément de base $(0, \ldots, 1, \ldots, 0)$ de
$\Gamma(Y, \mathcal{O}_{Y}^{n})$ ; cet homomorphisme définit évidemment un
homomorphisme de faisceaux
\[
  u : A^{n} \longrightarrow B .
\]

\note{cette formule est encadrée à la main sur le feuillet}

Les faisceaux $A^{n}$ et $B$ sont quasi-cohérents sur $Y$ (II, 1, 2, prop.~5),
donc cohérents, et il en est donc de même du noyau et du conoyau de $u$.

\marginal{puisque $Y$ est noethérien et $B$ de type fini par hypothèse}

Or, par construction, si \add{s est} une section de $\mathrm{Ker}\,u$ ou de
$\mathrm{Coker}\,u$ au-dessus d'un voisinage de $y$, le germe de $s$ au point
$y$ est nul, donc $s$ est nulle dans un voisinage \add{de} $y$ ; on en conclut
qu'il existe un voisinage ouvert $V$ de $y$ tel que la restriction $u|V$ soit un
\emph{isomorphisme} $A^{n}|V \to B|V$.

\note{\add{s est} est une correction manuscrite portée au-dessus de la ligne,
le tapuscrit ayant écrit \guillemotleft{} si une section \guillemotright}

Cela étant, en vertu du th.~2, il suffit de prouver que pour tout \ill{}

\marginal{\ill{}}

\note{une seconde note manuscrite, de deux mots, court obliquement dans la marge
gauche à mi-feuillet ; elle est au crayon, très pâle, et ne se laisse pas lire.
Une croix au crayon marque en outre le haut de la marge}

\page{4}

\note{demi-feuillet : le bas en est arraché. C'est le dictionnaire de la
dualité, écrit sans un mot de prose}

\struck{$X \subset$} $\mathbf{P}^{r} \supset X \supset Y$

\[
  H^{p}_{Y}(X, F) \simeq H^{p}_{Y}(P, F) = \varinjlim\;
  \mathrm{Ext}^{p}\!\left(P; \mathcal{O}_{Y_{n}}, F\right)
\]

son dual est
\[
  \left[\,H^{p}(X - Y, F) = \varinjlim\;
  \mathrm{Ext}^{p}\!\left(P; J^{n}, F\right)\,\right]
\]

\[
  \varprojlim\; \mathrm{Ext}^{r-p}\!\left(P; F,
  \mathcal{O}_{Y_{n}} \otimes \Omega^{r}_{P/k}\right)
  \;\xrightarrow{\ \sim\ }\;
  \mathrm{Ext}^{r-p}\!\left(\hat{Y}, \hat{F}, \hat{\Omega}^{r}_{P^{r}}\right)
\]

\[
  H^{p}_{Y}(X, F) \to H^{p}(X, F) \to H^{p}(X - Y, F) \to H^{p+1}_{Y}(X, F)
\]
\[
  \mathrm{Ext}^{r-p}_{\mathcal{O}_{P}}\!\left(\hat{Y}, \hat{F}, \hat{\Omega}^{r}\right)
  \leftarrow \mathrm{Ext}^{r-p}\!\left(P, F, \Omega^{r}\right)
  \leftarrow \varprojlim\; \mathrm{Ext}^{r-p}\!\left(P, F, J^{n}\Omega\right)
  \leftarrow
\]
\[
  \mathrm{Ext}^{r}
\]

\note{les deux dernières lignes sont écrites l'une sous l'autre, terme à terme :
la seconde est le dual de la première, et le feuillet s'arrête sur
$\mathrm{Ext}^{r}$ seul, là où la déchirure commence}

\page{5}

Premier encadré, deux suites exactes écrites l'une sous l'autre, appariées
terme à terme :

\begin{tikzcd}[column sep=small, row sep=small, nodes={font=\scriptsize}]
  \mathrm{Ext}^{i}(X, F, G) \arrow[r] \arrow[d, no head, "\times" description]
  & \mathrm{Ext}^{i}(X', F', G') \arrow[r] \arrow[d, no head, "\times" description]
  & \mathrm{Ext}^{i+1}(X, F, G) \arrow[d, no head, "\times" description] \\
  \mathrm{Ext}^{r-i}_{a}(X, G, F)
  & \mathrm{Ext}^{r-i}(X', G', F') \arrow[l]
  & \mathrm{Ext}^{r-1-i}(X, G, F) \arrow[l]
\end{tikzcd}

\note{les traits verticaux ne sont pas des flèches : ce sont les croix par
lesquelles il note l'accouplement de dualité entre les deux lignes. L'indice
$a$ des $\mathrm{Ext}$ de la seconde ligne est une petite lettre que le feuillet
ne définit nulle part ; elle est reproduite telle quelle}

\[
  \mathrm{Ext}^{i}_{Y}(X; F, G) = \varinjlim\; \mathrm{Ext}^{i}(X, F_{n}, G)
\]
dual à
\[
  \varprojlim\; \mathrm{Ext}^{r-i}_{a}(X, G, F_{n}) \simeq
  \mathrm{Ext}^{r-i}(\hat{X}, \hat{G}, \hat{F})
\]

\[
  \mathrm{Ext}^{i}_{Y'}(X'; F', G') = \varinjlim\; \mathrm{Ext}^{i}(X'; F'_{n}, G')
\]
dual à
\[
  \varprojlim\; \mathrm{Ext}^{r-1-i}(X', G', F'_{n})
\]
\[
  \mathrm{Ext}^{r-1-i}(\hat{X}', \hat{G}', \hat{F}') \;?
\]

\note{le point d'interrogation est de sa main et porte sur l'identification
elle-même : le passage à $X'$ est ce dont il n'est pas sûr}

Second encadré, de nouveau apparié :

\begin{tikzcd}[column sep=small, row sep=small, nodes={font=\scriptsize}]
  \mathrm{Ext}^{i}_{a}(X, F, G) \arrow[r]
  & \mathrm{Ext}^{i}_{Y}(X; F, G) \arrow[r] \arrow[d, no head, "\otimes" description]
  & \mathrm{Ext}^{i}_{Y'}(X'; F', G') \arrow[d, no head, "\otimes" description] \\
  \mathrm{Ext}^{r-i}(\hat{X}, \hat{G}, \hat{F})
  & \mathrm{Ext}^{r-i}_{a}(\hat{X}, \hat{G}, \hat{F}) \arrow[l]
  & \mathrm{Ext}^{r-1-i}(\hat{X}', \hat{G}', \hat{F}') \arrow[l]
\end{tikzcd}

\[
  \mathrm{Ext}^{r-i}(\hat{X}, \hat{G}, \hat{F}) \simeq \mathrm{Ext}^{r-i}(X; G, F)
\]

\struck{$\mathrm{Ext}^{r-i}(X; G, F)$}

\marginal{plus général}

\note{deux régions de cet encadré sont entourées au trait, sans que le feuillet
dise ce que l'entourage retient}

\[
  \mathrm{Ext}^{i}_{Y}(X; F, G) \longrightarrow \mathrm{Ext}^{i}_{Z}(X; F, G)
  \longrightarrow \mathrm{Ext}^{i}_{Z-Y}(X; F, G)
\]
\[
  \mathrm{Ext}^{r-1-i}_{a}\!\left(X_{/Y}; G_{/Y}, F_{/Y}\right) \longleftarrow
  \mathrm{Ext}^{r-1-i}_{a}\!\left(X_{/Z}; G_{/Z}, F_{/Z}\right) \longleftarrow
\]
\[
  X_{/Y} \longrightarrow X_{/Z}
\]

\page{6}

\note{second feuillet dactylographié, en tête \guillemotleft{} -- III-5 -- \guillemotright,
annoté et souligné au feutre comme le précédent. Ce qu'il prouve — que la
cohomologie du complémentaire de $V(f)$ se calcule par le complexe de Koszul de
la famille $f$ — est exactement la machine dont les feuillets manuscrits se
servent, ce qui est la raison pour laquelle ces deux feuillets sont transcrits et
non passés}

$C^{p}(\mathfrak{U}, F)$ étant le groupe des $p$-cochaînes \emph{alternées} de
Čech, relativement au recouvrement $\mathfrak{U}$ et à coefficients dans $F$,
il résulte de ce qui précède que l'on peut écrire
\[
  (11) \qquad C^{p}(\mathfrak{U}, F) = \varinjlim_{n}\; C^{p}_{n}(M) .
\]
Or, avec les notations du n\textsuperscript{o}~1, $C^{p}_{n}(M)$ s'identifie à
$K^{p+1}(f, M)$, et l'application $\varphi_{mn}$ s'identifie à l'application
$\varphi^{n-m}$ définie au n\textsuperscript{o}~1. On a donc, pour tout
$p \geqslant 0$, un isomorphisme canonique fonctoriel en $M$
\[
  (12) \qquad C^{p}(\mathfrak{U}, F) \;\xrightarrow{\ \sim\ }\;
  C^{p+1}\!\left((f), M\right) .
\]
La formule (3) du n\textsuperscript{o}~1 montre en outre que ces isomorphismes
sont compatibles avec les opérateurs cobord. Il en résulte :

\textbf{Proposition 2.} Si $X$ est \add{un préschéma} noethérien, ou
\struck{\ill{}} un schéma dont l'espace de base est quasi-compact, il existe
pour tout $p \geqslant 1$ un isomorphisme canonique fonctoriel en $F$
\[
  (13) \qquad H^{p}(\mathfrak{U}, F) \;\xrightarrow{\ \sim\ }\;
  H^{p+1}\!\left((f), M\right) .
\]
On a en outre une suite exacte fonctorielle en $F$
\[
  (14) \qquad 0 \to H^{0}\!\left((f), M\right) \to M \to
  H^{0}(\mathfrak{U}, F) \to H^{1}\!\left((f), M\right) \to 0 .
\]

\note{\add{un préschéma} est une addition manuscrite au-dessus de la ligne ;
elle corrige \guillemotleft{} si $X$ est noethérien \guillemotright{} en
\guillemotleft{} si $X$ est un préschéma noethérien \guillemotright, et un
membre de phrase tapé est biffé à la machine juste après}

Les relations (13) sont en effet évidentes d'après ce qui précède. D'autre part,
\struck{\ill{}} $M = C^{0}((f), M)$ et $C^{0}(\mathfrak{U}, F) = C^{1}((f), M)$,
donc $H^{0}(\mathfrak{U}, F)$ s'identifie au sous-groupe des $1$-cocycles de
$C^{1}((f), M)$ ; la suite exacte (3) n'est autre alors que
\add{celle qui résulte de} la définition des groupes de cohomologie
$H^{0}((f), M)$ et $H^{1}((f), M)$.

\textbf{Corollaire.} Supposons qu'il existe des
$g_{i} \in \Gamma(\mathfrak{U}, \mathcal{O}_{X})$ tels que
$\sum_{i=1}^{n} g_{i} f_{i} = 1$ \add{au-dessus de $\mathfrak{U}$}. Alors, pour
tout faisceau quasi-cohérent $G$ sur $\mathfrak{U}$, on a
$H^{i}(\mathfrak{U}, G) = 0$ pour $i > 0$ ; si en outre $X = \mathfrak{U}$,
l'homomorphisme canonique $M \to H^{0}(\mathfrak{U}, F)$ est bijectif.
\add{(les $X_{f_{i}} = \mathfrak{U}_{i}$ étant par hypothèse quasi-compacts)}

On peut en effet supposer $X = \mathfrak{U}$, et alors $H^{i}((f), M) = 0$ pour
tout $i \geqslant 0$ (n\textsuperscript{o}~1, cor.\ de la prop.~1) ; l'assertion
découle alors aussitôt de (13) et (14). On observera que \add{puisque}
$H^{0}(\mathfrak{U}, F) = H^{0}(\mathfrak{U}, F)$, on redémontre ainsi le \ill{}

\note{les additions \add{au-dessus de $\mathfrak{U}$}, \add{celle qui résulte de},
\add{(pacts)} — abrégeant \emph{quasi-compacts} — et \add{puisque} sont toutes
manuscrites, portées au-dessus de la ligne ou en marge}

\page{7}

\struck{$\mathrm{Ext}^{r-i}\!\left(\hat{P}, \Omega^{r}_{X_{P}}(-n)\right)$}

\[
  H^{r-i}\!\left(\hat{P}, \Omega^{r}_{P}(-n)\right)
\]

\struck{$H^{p}(\hat{P},$} \quad \struck{$H^{p}(\hat{P},$}

Encadré :
\[
  H^{i}\!\left(\hat{P}, \Omega^{r}_{P}(-n)\right) = 0 \quad \text{si } 0 < i < r-m
\]
\[
  H^{r}\!\left(P, \Omega^{r}(-n)\right) \;\xrightarrow{\ \sim\ }\;
  H^{r}\!\left(\hat{P}, \hat{\Omega}^{r}(-n)\right)
\]
pour $n$ grand, \ill{} et \ill{} (en somme) les \ill{} \ill{} remplaçant
$\Omega^{r}_{P}$ par $\mathcal{O}_{P}$ \ill{}

\marginal{$\Longleftrightarrow$ $P - X$ est de dim.\ coh.\ $r-m$
[\,et \ill{} pour \ill{} grand\,]}

Les schémas formels s'introduisent en \uncertain{géométrie}.
\uncertain{Je les groupe} de la façon suivante :

\begin{enumerate}
  \item[1)] Techniquement, pour énoncer et prouver des th.\ généraux comme les
    deux th.\ fondamentaux des morphismes propres.
  \item[3)] En \uncertain{théorie formelle} des Modules, \ill{}
    \struck{\ill{}} une notion de familles \ill{}
  \item[2)] En théorie des Lefschetz -- \uncertain{Grauert}, \ill{}
    \add{le plus \uncertain{couramment}} \ill{} \emph{relations} entre une
    variété (projective, ou locale) et son complété formel le long d'une
    \guillemotleft{} section hyperplane \guillemotright.
  \item[4)] En formalisme de dualité, les énoncés où les groupes de cohomologie
    $H^{p}_{Y}(X, F)$ s'expriment par des $\mathrm{Ext}$ ou $H^{p}$ sur
    $\hat{X}$ (complété formel le long de $Y$).
\end{enumerate}

\note{l'ordre est celui du feuillet, et il n'est pas croissant : l'item 3 est
écrit avant l'item 2, une flèche partant de la marge gauche le rappelle vers le
haut, et le \emph{2} de l'item 2 est entouré. Le nom qui suit
\emph{Lefschetz} est écrit d'une main rapide ; les lettres sont compatibles
avec \emph{Grauert}, que la matière rendrait naturel, mais le feuillet ne le
prouve pas}

\page{8}

\begin{enumerate}
  \item[5)] Théorie des groupes formels, où \uncertain{les} th.\ \ill{} la
    \ill{} sont \ill{} des schémas formels \ldots
\end{enumerate}

De plus, notons

\begin{enumerate}
  \item[6)] Dans la théorie des types de \emph{Greenberg} -- \emph{Néron}
    s'introduisent des schémas relatifs sur des \ill{} anneaux, qui ne sont pas
    nécessairement \ill{} des précédents \ldots
\end{enumerate}

\note{\emph{Greenberg} et \emph{Néron} sont soulignés de sa main. Le feuillet
s'arrête là : les deux tiers inférieurs en sont blancs, et le programme des six
items n'a pas de suite dans le dossier}

\end{document}
